\section{Data contracts and observable quantities}
\label{sec:data-contracts}

The contract catalog separates fourteen boundaries because similarly shaped arrays can represent different experiments. Its semantic version belongs to the interface; a payload's schema tag identifies its serialized family. These identifiers need not have the same numeric suffix. Normative JSON Schemas validate committed JSON and generate TypeScript types. A generated type is not a runtime validator for an arbitrary network response. Numerical reconciliation, subject ownership, temporal eligibility, graph identity and licensing are separate checks.

The tables below are generated from the catalog and schema files by \path{scripts/build_report_contract_tables.py}. They list every root field, its declared type, root constraint, required status and quantity category. Nested arrays and objects follow the corresponding normative schema; the prose states the scientific meaning that a structural type cannot supply. $U$ denotes a source series' measurement unit, which remains unspecified when the admitted provenance does not identify a physical unit. Counts and observation positions are distinguished from measured values. Statistical ratios and normalized residuals are dimensionless.

\subsection{Trace: displayed history and retained predictions}
\label{sec:contract-trace}
The trace boundary records one selected case's display history, reserved targets, method forecasts and summaries. Table~\ref{tab:contract-trace} describes interface revision \texttt{2.0.0}. Its history length is the length of the original prefix, while the exported history may contain fewer display samples. The parallel history-index array supplies original observation positions. Those positions are essential: a three-hundred-point display is not automatically a three-hundred-step contiguous fitting sequence. Forecast paths and targets retain source units; seasonal lag and horizon are observation counts, and quantile levels are probabilities. Method records combine prediction arrays with scaled and unscaled scores, so the container has mixed units rather than one inferred physical unit. Redaction is explicit and does not license a consumer to reconstruct or expose restricted source values. The pipeline's trace serialization is checked against the schema and cross-field tests before browser copying. The atlas uses the admitted fresh comparison for contiguous fitting history and verifies its target relationship to the trace. A missing case does not acquire another trace's identity, and a retrospective diagnostic result is not silently substituted for a forecast path.
\begingroup\small
\begin{longtable}{@{}p{.22\textwidth}p{.14\textwidth}p{.14\textwidth}p{.29\textwidth}p{.09\textwidth}@{}}
\caption{Root fields of the trace contract, interface revision \texttt{2.0.0}. Types, units and structural constraints distinguish required members from optional metadata; nested records are listed separately.}\label{tab:contract-trace}\\
\toprule Field & JSON type & Unit & Root constraint & Required\\\midrule\endfirsthead
\toprule Field & JSON type & Unit & Root constraint & Required\\\midrule\endhead
\texttt{schema} & fixed value & metadata & const="chronosco\allowbreak{}pe.trace/v2" & yes\\
\texttt{case\_\allowbreak{}id} & string & metadata & No explicit range & yes\\
\texttt{seasonality} & integer & observations & minimum=1 & yes\\
\texttt{horizon} & integer & observations & minimum=1 & yes\\
\texttt{quantile\_\allowbreak{}levels} & array & probabilities &  & yes\\
\texttt{redacted} & boolean & metadata & No explicit range & yes\\
\texttt{history\_\allowbreak{}len} & integer & observations & minimum=0 & yes\\
\texttt{history\_\allowbreak{}index} & array & indices &  & yes\\
\texttt{history} & array & source U &  & yes\\
\texttt{actual} & array & source U &  & yes\\
\texttt{methods} & array & mixed; see text &  & yes\\
\texttt{summary} & object reference & mixed; see text & No explicit range & yes\\
\bottomrule
\end{longtable}
\endgroup


\subsection{Manifest: subject identity, provenance and delivery references}
\label{sec:contract-manifest}
The manifest binds a case identifier to source classification, provenance, engine identity, series descriptors, artifact references and an execution-lane decision. Table~\ref{tab:contract-manifest} describes interface revision \texttt{2.0.0}. Its best-method and metric fields summarize an existing experiment; they are not permission to train a selector using the reserved block. Series descriptors have mixed units: counts describe lengths, distribution summaries inherit source units, and correlations are dimensionless. The seed is a reproducibility identifier rather than a physical measurement. Artifact references include paths and byte counts, with analysis and streaming references allowed to represent absence under their declared forms. A public-artifact permission belongs to provenance and must be checked separately from whether a JSON object parses. The browser loads the selected manifest before its dependent trace and checks that the returned subject matches the requested admitted case. Index membership prevents an arbitrary case path from becoming a selected subject. Missing optional analysis differs from a missing core trace. Producer tests protect the manifest and reference structure, while consumer tests protect case consistency; neither a human-readable category nor an engine label independently establishes that two payloads share the same numerical inputs.
\begingroup\small
\begin{longtable}{@{}p{.22\textwidth}p{.14\textwidth}p{.14\textwidth}p{.29\textwidth}p{.09\textwidth}@{}}
\caption{Root fields of the manifest contract, interface revision \texttt{2.0.0}. Types, units and structural constraints distinguish required members from optional metadata; nested records are listed separately.}\label{tab:contract-manifest}\\
\toprule Field & JSON type & Unit & Root constraint & Required\\\midrule\endfirsthead
\toprule Field & JSON type & Unit & Root constraint & Required\\\midrule\endhead
\texttt{schema} & fixed value & metadata & const="chronosco\allowbreak{}pe.manifest/v1" & yes\\
\texttt{case\_\allowbreak{}id} & string & metadata & No explicit range & yes\\
\texttt{category} & string & metadata & No explicit range & yes\\
\texttt{real\_\allowbreak{}or\_\allowbreak{}synthetic} & string & metadata & No explicit range & yes\\
\texttt{expected\_\allowbreak{}band} & string & metadata & No explicit range & yes\\
\texttt{engine} & object & metadata &  & yes\\
\texttt{provenance} & object & metadata &  & yes\\
\texttt{analysis\_\allowbreak{}artifact} & declared union & metadata & No explicit range & yes\\
\texttt{streaming\_\allowbreak{}artifact} & declared union & metadata & No explicit range & yes\\
\texttt{series} & object reference & mixed; see text & No explicit range & yes\\
\texttt{seed} & integer & metadata & minimum=0 & yes\\
\texttt{artifact} & object reference & metadata & No explicit range & yes\\
\texttt{lane} & fixed value & metadata & enum=["live","pr\allowbreak{}ecompute"] & yes\\
\texttt{gate} & object reference & metadata & No explicit range & yes\\
\texttt{flags} & array & metadata &  & yes\\
\texttt{best\_\allowbreak{}method} & string/null & metadata & No explicit range & yes\\
\texttt{metrics} & object & mixed; see text &  & yes\\
\bottomrule
\end{longtable}
\endgroup


\subsection{Case index: inventory rather than an accuracy claim}
\label{sec:contract-index}
The case index is the compact entry point for browsing the selected corpus. Table~\ref{tab:contract-index} describes interface revision \texttt{1.0.0}. It contains a schema tag, engine-version metadata, a declared case count and the case entries. The count is a number of admitted subjects, not a number of independent forecasting trials. Entry descriptors support grouping and navigation; scientific comparison still requires the associated trace and experiment records. The index is produced from the committed case inventory and consumed by the source selector. A loading failure leaves the inventory unavailable rather than creating a plausible set of cases. An unknown requested identifier is refused before a manifest request. This matters for shared links: displaying the first available case under a stale requested identifier would give the wrong subject to otherwise valid forecasts. Structural validation can check field types and required members, but reconciliation must also compare the declared count with actual entries and referenced manifests. The original fifteen cases remain separate from the external combination inventory and the complete daily corpus. A larger count in either of those experiments does not change what this selected-case index represents.
\begingroup\small
\begin{longtable}{@{}p{.22\textwidth}p{.14\textwidth}p{.14\textwidth}p{.29\textwidth}p{.09\textwidth}@{}}
\caption{Root fields of the index contract, catalog version 1.0.0. Types, units and structural constraints distinguish required members from optional metadata; nested records follow the cited schema.}\label{tab:contract-index}\\
\toprule Field & JSON type & Unit & Root constraint & Required\\\midrule\endfirsthead
\toprule Field & JSON type & Unit & Root constraint & Required\\\midrule\endhead
\texttt{schema} & fixed value & metadata & const="chronosco\allowbreak{}pe.index/v1" & yes\\
\texttt{engine\_\allowbreak{}version} & string & metadata & No additional root-field range & yes\\
\texttt{n\_\allowbreak{}cases} & integer & counts & minimum=0 & yes\\
\texttt{cases} & array & metadata & Nested content follows the normative schema & yes\\
\bottomrule
\end{longtable}
\endgroup


\subsection{Analysis: heterogeneous diagnostics and explicit availability}
\label{sec:contract-analysis}
The analysis boundary holds retrospective diagnostics for one supplied sequence. Table~\ref{tab:contract-analysis} describes interface revision \texttt{2.1.0}. Its required groups cover stationarity, autocorrelation, seasonality, filtering, change points, volatility, distribution, fractal measures, nonlinear analysis and causality. An optional sequence group carries the later original-axis pattern calculations. These objects cannot share a single numerical unit: a spectral frequency, correlation coefficient, residual variance, entropy estimate and raw time-reversal statistic have different dimensions. The sequence length and declared seasonal period are observation counts. Method descriptions define the remaining quantities, availability rules and assumptions. Some groups are deliberately flexible objects, so a structurally valid analysis file is not a certificate that every numerical diagnostic is meaningful on every input. Missing dependencies and insufficient samples must remain visible as unavailable calculations. The analysis producer runs outside the browser for heavier methods; lightweight browser diagnostics are independent computations with their own scope. Several stored analyses describe the complete sequence, including later forecast targets. They are therefore not admissible as origin-time predictors without a separate prefix-only computation. The consumer never turns such retrospective information into evidence that a forecast issued earlier knew the future.
\begingroup\small
\begin{longtable}{@{}p{.22\textwidth}p{.14\textwidth}p{.14\textwidth}p{.29\textwidth}p{.09\textwidth}@{}}
\caption{Root fields of the analysis contract, catalog version 2.1.0. Types, units and structural constraints distinguish required members from optional metadata; nested records follow the cited schema.}\label{tab:contract-analysis}\\
\toprule Field & JSON type & Unit & Root constraint & Required\\\midrule\endfirsthead
\toprule Field & JSON type & Unit & Root constraint & Required\\\midrule\endhead
\texttt{n} & integer & observations & minimum=1 & yes\\
\texttt{seasonality\_\allowbreak{}period} & integer & observations & minimum=1 & yes\\
\texttt{stationarity} & object & metadata & Nested content follows the normative schema & yes\\
\texttt{autocorrelation} & object & metadata & Nested content follows the normative schema & yes\\
\texttt{seasonality} & object & observations & Nested content follows the normative schema & yes\\
\texttt{filters} & object & metadata & Nested content follows the normative schema & yes\\
\texttt{changepoints} & object & metadata & Nested content follows the normative schema & yes\\
\texttt{volatility} & object & metadata & Nested content follows the normative schema & yes\\
\texttt{distribution} & object & metadata & Nested content follows the normative schema & yes\\
\texttt{fractal} & object & metadata & Nested content follows the normative schema & yes\\
\texttt{nonlinear} & object & metadata & Nested content follows the normative schema & yes\\
\texttt{causality} & object & metadata & Nested content follows the normative schema & yes\\
\texttt{sequence} & declared union & metadata & No additional root-field range & no\\
\bottomrule
\end{longtable}
\endgroup


\subsection{Streaming: aggregate trajectories from a specific replay}
\label{sec:contract-streaming}
The streaming boundary records nominal coverage, horizon, warmup, method trajectories, optional covariate information and reference metadata. Table~\ref{tab:contract-streaming} describes interface revision \texttt{2.0.0}. Coverage is a fraction; the horizon and warmup count observations. Method objects retain metric trajectories and timing-related quantities whose units belong to their named measures. The optional covariate distinguishes absence from an aligned driver and must not imply that future values were known unless its arrival policy says so. This serialized family represents the historical replay configuration, including the smoothing-plus-half-drift base and the evaluated adaptive controllers. It is not interchangeable with the corrected finite-history Theta forecaster in the common final-block experiment. Nor does a stored latency establish a resource-normalized comparison without its execution environment. The producer's schema and numerical tests protect the historical artifact; the later per-step reconstruction supplies a separate detailed audit. If streaming output is missing, the web view reports unavailability rather than simulating replacement trajectories. The nominal target is not a mathematical guarantee for each finite window, and temporal averages must retain their distinction from averages of local absolute coverage errors.
\begingroup\small
\begin{longtable}{@{}p{.22\textwidth}p{.14\textwidth}p{.14\textwidth}p{.29\textwidth}p{.09\textwidth}@{}}
\caption{Root fields of the streaming contract, interface revision \texttt{2.0.0}. Types, units and structural constraints distinguish required members from optional metadata; nested records are listed separately.}\label{tab:contract-streaming}\\
\toprule Field & JSON type & Unit & Root constraint & Required\\\midrule\endfirsthead
\toprule Field & JSON type & Unit & Root constraint & Required\\\midrule\endhead
\texttt{nominal\_\allowbreak{}coverage} & number & probabilities & minimum=0; maximum=1 & yes\\
\texttt{horizon} & integer & observations & minimum=1 & yes\\
\texttt{warmup} & integer & metadata & minimum=0 & yes\\
\texttt{methods} & object & mixed; see text &  & yes\\
\texttt{covariate} & object/null & mixed; see text & Open object & yes\\
\texttt{references} & object & metadata &  & yes\\
\bottomrule
\end{longtable}
\endgroup


\subsection{Browser point-model metadata and binary execution}
\label{sec:contract-nlinear-model}
The browser point-model boundary pairs legacy metadata with an ONNX binary. Table~\ref{tab:contract-nlinear-model} describes interface revision \texttt{2.0.0}. The metadata retains its historical name while the interface identifies the computation as a normalized linear map. It declares lookback, horizon, normalization, normalized residual scale and training-window count. Counts and the residual scale are dimensionless, with lookback and horizon measured in observations. The historical metadata schema permits positive dimensions generally, but the admitted graph and current calibration have the specific forty-eight-input, twelve-output identity. The graph is a binary asset; JSON validation of its companion metadata is not ONNX execution or numerical parity. The browser verifies the graph digest through the calibration binding and checks actual output dimensions, type and finite values. The original point matrix remains unchanged by the new interval study. A missing model is omitted from the available live paths; it does not replace another method or acquire a favorable score. Constant lookbacks and unsupported horizons are refused because the required normalization or output contract is unavailable. Legacy Gaussian dispersion is retained as historical metadata, not reused as a silent fallback for the confirmed interval procedure.
\begingroup\small
\begin{longtable}{@{}p{.22\textwidth}p{.14\textwidth}p{.14\textwidth}p{.29\textwidth}p{.09\textwidth}@{}}
\caption{Root fields of the nlinear-model contract, interface revision \texttt{2.0.0}. Types, units and structural constraints distinguish required members from optional metadata; nested records are listed separately.}\label{tab:contract-nlinear-model}\\
\toprule Field & JSON type & Unit & Root constraint & Required\\\midrule\endfirsthead
\toprule Field & JSON type & Unit & Root constraint & Required\\\midrule\endhead
\texttt{name} & fixed value & metadata & const="NLinear" & yes\\
\texttt{lookback} & integer & observations & minimum=1 & yes\\
\texttt{horizon} & integer & observations & minimum=1 & yes\\
\texttt{normalization} & fixed value & metadata & const="subtract\_\allowbreak{}last\_\allowbreak{}then\_\allowbreak{}divide\_\allowbreak{}std" & yes\\
\texttt{sigma\_\allowbreak{}norm} & number & dimensionless & minimum=0 & yes\\
\texttt{trained\_\allowbreak{}windows} & integer & counts & minimum=1 & yes\\
\bottomrule
\end{longtable}
\endgroup


\subsection{Confirmed browser calibration: dimensions and statistical population}
\label{sec:contract-browser-calibration}
The confirmed-calibration boundary binds an interval procedure to one exact point graph. Table~\ref{tab:contract-browser-calibration} describes interface revision \texttt{1.0.0}. Required fields fix the graph digest, observed-window normalization, lookback, horizon, miscoverage, roughness thresholds, calibration counts and absolute-residual radii. Radii are normalized by the observed-prefix standard deviation; multiplying by that scale restores source units. Roughness is dimensionless because both its numerator and denominator have squared source units. Three groups each carry twelve lead-specific radii, while global radii retain a separate comparator. The producer uses disjoint calibration and test identities, and independent evidence checks reconstruct all radii and held-out paths. The consumer refuses a mismatched graph or malformed calibration instead of returning to Gaussian bands. This structural binding prevents accidental model-calibration substitution, but it does not prove that an authentic deployment series is exchangeable with the generated calibration population. Overall coverage, family outcomes and lead outcomes remain statistical evidence with that population boundary. A fixed probability level in the schema also does not authorize interpolation to another level or a simultaneous whole-path coverage claim.
\begingroup\small
\begin{longtable}{@{}p{.22\textwidth}p{.14\textwidth}p{.14\textwidth}p{.29\textwidth}p{.09\textwidth}@{}}
\caption{Root fields of the browser-calibration contract, catalog version 1.0.0. Types, units and structural constraints distinguish required members from optional metadata; nested records follow the cited schema.}\label{tab:contract-browser-calibration}\\
\toprule Field & JSON type & Unit & Root constraint & Required\\\midrule\endfirsthead
\toprule Field & JSON type & Unit & Root constraint & Required\\\midrule\endhead
\texttt{schema} & fixed value & metadata & const="chronosco\allowbreak{}pe.browser-calib\allowbreak{}ration/v1" & yes\\
\texttt{point\_\allowbreak{}model\_\allowbreak{}sha256} & fixed value & metadata & const="621265165\allowbreak{}211f756cff537ade\allowbreak{}e364b9790c34ced5\allowbreak{}67bc2a5087513360\allowbreak{}e1f69f6" & yes\\
\texttt{lookback} & fixed value & observations & const=48 & yes\\
\texttt{horizon} & fixed value & observations & const=12 & yes\\
\texttt{alpha} & fixed value & probabilities & const=0.2 & yes\\
\texttt{normalization} & fixed value & metadata & const="subtract\_\allowbreak{}last\_\allowbreak{}then\_\allowbreak{}divide\_\allowbreak{}population\_\allowbreak{}std" & yes\\
\texttt{roughness\_\allowbreak{}boundaries} & fixed value & dimensionless & const=[0.5,1.25] & yes\\
\texttt{calibration\_\allowbreak{}bin\_\allowbreak{}counts} & array & counts & minItems=3; maxItems=3 & yes\\
\texttt{absolute\_\allowbreak{}residual\_\allowbreak{}radii} & array & dimensionless & minItems=3; maxItems=3 & yes\\
\texttt{global\_\allowbreak{}radii} & array & dimensionless & minItems=12; maxItems=12 & yes\\
\bottomrule
\end{longtable}
\endgroup


\subsection{Ridge reference study: a retained unsuccessful point candidate}
\label{sec:contract-browser-ridge-study}
The ridge-study boundary preserves the independent point-model selection experiment even though its selected point map did not improve the declared criterion. Table~\ref{tab:contract-browser-ridge-study} describes interface revision \texttt{1.1.0}. Its JSON identifies the candidate, fixed input and output dimensions, miscoverage, lead-specific absolute-residual radii and seven selection records. The optional explicit normalization descriptor documents the producer's actual convention. A companion ONNX binary makes the selected float32 map executable without implying that it is the served point model. Selection losses and residual radii are normalized quantities; horizon and lookback count observations. Training, validation, calibration and test roles remain distinct in the numerical evidence. The publisher and independent checker bind source hashes, split identities, graph structure and scored forecasts. Removing an unfavorable candidate or replacing its stored outcomes with a later successful interval experiment would erase the question this contract exists to preserve. The browser's production calibration is a different boundary with a different point-graph identity. Missing or changed ridge evidence is therefore an admission failure for the study, not permission to substitute the live graph and continue reporting the original candidate's results.
\begingroup\small
\begin{longtable}{@{}p{.22\textwidth}p{.14\textwidth}p{.14\textwidth}p{.29\textwidth}p{.09\textwidth}@{}}
\caption{Root fields of the browser-ridge-study contract, catalog version 1.1.0. Types, units and structural constraints distinguish required members from optional metadata; nested records follow the cited schema.}\label{tab:contract-browser-ridge-study}\\
\toprule Field & JSON type & Unit & Root constraint & Required\\\midrule\endfirsthead
\toprule Field & JSON type & Unit & Root constraint & Required\\\midrule\endhead
\texttt{name} & fixed value & metadata & const="Held-out normalized ridge" & yes\\
\texttt{lookback} & fixed value & observations & const=48 & yes\\
\texttt{horizon} & fixed value & observations & const=12 & yes\\
\texttt{alpha} & fixed value & probabilities & const=0.2 & yes\\
\texttt{absolute\_\allowbreak{}residual\_\allowbreak{}radii} & array & dimensionless & minItems=12; maxItems=12 & yes\\
\texttt{selection} & array & metadata & minItems=7; maxItems=7 & yes\\
\texttt{normalization} & fixed value & metadata & const="subtract\_\allowbreak{}last\_\allowbreak{}then\_\allowbreak{}divide\_\allowbreak{}population\_\allowbreak{}std" & no\\
\bottomrule
\end{longtable}
\endgroup


\subsection{Archived comparison: common targets with mixed execution provenance}
\label{sec:contract-reevaluation}
The archived comparison boundary represents a common final-block rescoring experiment. Table~\ref{tab:contract-reevaluation} describes interface revision \texttt{2.1.0} and payload family v2. It records the protocol, source engine version, trace and implementation digests, full case records, summaries and input identities. Reconstructed contiguous histories and common targets distinguish it from a comparison of incompatible rolling-backtest aggregates. Targets and prediction paths have source units; MAE and RMSE retain those units, while MASE, normalized quantile loss and scaled interval score are dimensionless. Empirical coverage is a fraction. Heavy paths in this experiment remain frozen, while specified lightweight paths are recomputed. Their coexistence is explicit provenance, not a claim that every model was freshly executed. Independent numerical checks reconcile stored paths, historical scales, targets and summaries. The web consumer labels the archived experiment separately from the fresh one and does not silently use it when fresh results are unavailable. Optional later audit fields extend the schema without changing the meaning of the original calculation. A source hash binds bytes, but does not by itself establish representative sampling, absence of pretraining overlap or a justified population uncertainty interval.
\begingroup\small
\begin{longtable}{@{}p{.22\textwidth}p{.14\textwidth}p{.14\textwidth}p{.29\textwidth}p{.09\textwidth}@{}}
\caption{Root fields of the reevaluation contract, interface revision \texttt{2.1.0}. Types, units and structural constraints distinguish required members from optional metadata; nested records are listed separately.}\label{tab:contract-reevaluation}\\
\toprule Field & JSON type & Unit & Root constraint & Required\\\midrule\endfirsthead
\toprule Field & JSON type & Unit & Root constraint & Required\\\midrule\endhead
\texttt{schema} & fixed value & metadata & const="chronosco\allowbreak{}pe.reevaluation/\allowbreak{}v2" & yes\\
\texttt{protocol} & object & metadata &  & yes\\
\texttt{source\_\allowbreak{}engine\_\allowbreak{}version} & string & metadata & No explicit range & yes\\
\texttt{trace\_\allowbreak{}sha256} & object & metadata &  & yes\\
\texttt{code\_\allowbreak{}sha256} & object & metadata &  & yes\\
\texttt{cases} & array & metadata & minItems=1 & yes\\
\texttt{summary} & array & mixed; see text & minItems=1 & yes\\
\texttt{input\_\allowbreak{}sha256} & object & metadata &  & yes\\
\texttt{source\_\allowbreak{}comparison\_\allowbreak{}sha256} & string & metadata & pattern="\textasciicircum{}[0-9a-f]\{64\}\$" & no\\
\texttt{runtime} & object & metadata &  & no\\
\bottomrule
\end{longtable}
\endgroup


\subsection{Fresh comparison: execution identity as an additional invariant}
\label{sec:contract-fresh-reevaluation}
The fresh-comparison catalog row has interface revision \texttt{1.0.0} and intentionally uses the comparison schema rather than inventing another incompatible case format. Table~\ref{tab:contract-fresh-reevaluation} shows that shared structure. Its scientific distinction is the producer and admission contract: every configured variant is fitted or inferred on the stated original history, with package, checkpoint, source and runtime records retained. Per-method fitting and inference duration is reported in seconds for the specified execution environment; it is not an environment-independent computational complexity measurement. Exact identities and the complete inventory protect the difference between a new execution and an archived path with renamed provenance. The source provider additionally verifies case identity, horizon, seasonal lag, history length and target agreement before exposing the comparison under a selected subject. The fresh target values retain their own precision for scoring and are checked against the released trace's display precision. Missing, partial or stale fresh evidence fails admission; an archived result remains an explicitly selected separate experiment. Matching aggregate values can confirm numerical recomputation while still representing the same source series and reserved outcomes, so they are not independent statistical replications.
\begingroup\small
\begin{longtable}{@{}p{.22\textwidth}p{.14\textwidth}p{.14\textwidth}p{.29\textwidth}p{.09\textwidth}@{}}
\caption{Root fields of the fresh-reevaluation contract, catalog version 1.0.0. Types, units and structural constraints distinguish required members from optional metadata; nested records follow the cited schema.}\label{tab:contract-fresh-reevaluation}\\
\toprule Field & JSON type & Unit & Root constraint & Required\\\midrule\endfirsthead
\toprule Field & JSON type & Unit & Root constraint & Required\\\midrule\endhead
\texttt{schema} & fixed value & metadata & const="chronosco\allowbreak{}pe.reevaluation/\allowbreak{}v2" & yes\\
\texttt{protocol} & object & metadata & Nested content follows the normative schema & yes\\
\texttt{source\_\allowbreak{}engine\_\allowbreak{}version} & string & metadata & No additional root-field range & yes\\
\texttt{trace\_\allowbreak{}sha256} & object & metadata & Nested content follows the normative schema & yes\\
\texttt{code\_\allowbreak{}sha256} & object & metadata & Nested content follows the normative schema & yes\\
\texttt{cases} & array & metadata & minItems=1 & yes\\
\texttt{summary} & array & mixed; see text & minItems=1 & yes\\
\texttt{input\_\allowbreak{}sha256} & object & metadata & Nested content follows the normative schema & yes\\
\texttt{source\_\allowbreak{}comparison\_\allowbreak{}sha256} & string & metadata & pattern="\textasciicircum{}[0-9a-f]\{64\}\$" & no\\
\texttt{runtime} & object & metadata & Nested content follows the normative schema & no\\
\bottomrule
\end{longtable}
\endgroup


\subsection{Historical selector: training metadata rather than an oracle}
\label{sec:contract-selector}
The historical selector boundary records a method-ranking artifact produced by an offline training calculation. Table~\ref{tab:contract-selector} describes interface revision \texttt{1.0.0}. Its kind, metric and scope fields are fixed identifiers; ranking pairs contain method names and dimensionless scaled errors. The default method is a selected identifier, and the training-case count is a count rather than a performance confidence measure. This legacy artifact is not used to assign the winner of the common reserved block. In particular, a future-informed minimum over methods is an oracle summary and must not be presented as a deployable selector merely because both outputs contain method names. The later daily-corpus selector has its own discovery and confirmation protocol in a different artifact. Keeping the boundaries separate prevents an old default ranking from acquiring that newer experiment's validation claim. A missing selector does not change the admitted case comparison. Rebuilding it is an explicit offline operation, not something startup or deployment silently performs. The producer's contract checks fixed metadata and ranking structure; meaningful generalization still depends on which cases were available for training and which outcomes remained reserved for evaluating the selection rule.
\begingroup\small
\begin{longtable}{@{}p{.22\textwidth}p{.14\textwidth}p{.14\textwidth}p{.29\textwidth}p{.09\textwidth}@{}}
\caption{Root fields of the selector contract, catalog version 1.0.0. Types, units and structural constraints distinguish required members from optional metadata; nested records follow the cited schema.}\label{tab:contract-selector}\\
\toprule Field & JSON type & Unit & Root constraint & Required\\\midrule\endfirsthead
\toprule Field & JSON type & Unit & Root constraint & Required\\\midrule\endhead
\texttt{kind} & fixed value & metadata & const="method-selector" & yes\\
\texttt{metric} & fixed value & metadata & const="mean\_\allowbreak{}mase\_\allowbreak{}over\_\allowbreak{}training\_\allowbreak{}cases" & yes\\
\texttt{scope} & fixed value & metadata & const="classical" & yes\\
\texttt{ranking} & array & mixed; see text & Nested content follows the normative schema & yes\\
\texttt{default\_\allowbreak{}method} & string & metadata & No additional root-field range & yes\\
\texttt{n\_\allowbreak{}train\_\allowbreak{}cases} & integer & counts & minimum=1 & yes\\
\bottomrule
\end{longtable}
\endgroup


\subsection{Daily corpus: complete accounting and partition identity}
\label{sec:contract-forecastability}
The daily-corpus boundary holds the complete source accounting, protocol, implementation and package identities, excluded records, model roster, summaries, per-series rows and audit metadata. Table~\ref{tab:contract-forecastability} describes interface revision \texttt{1.0.0}. Its source count is fixed at 4,227 identities; completeness is a required state rather than an inference from a large displayed table. Eligibility, duplicate handling and undefined scales can change the included denominator, so excluded identities remain part of the result. Forty prefix features have heterogeneous definitions and availability conditions. Association targets and MASE summaries are dimensionless, while feature-specific units are stated by their mathematical construction. Discovery and confirmation assignments bind exact sequence identities. Multiple-testing adjustment is performed separately for the declared feature family in each partition. Those corrected associations do not automatically imply useful method selection on new series. The producer and independent checker reconstruct the sufficient statistics and retain the unfavorable confirmation-selector result. The browser consumes the complete corpus artifact rather than substituting the fifteen selected examples. A missing or incomplete file is unavailable evidence, not an opportunity to infer the missing population from a convenient sample.
\begingroup\small
\begin{longtable}{@{}p{.22\textwidth}p{.14\textwidth}p{.14\textwidth}p{.29\textwidth}p{.09\textwidth}@{}}
\caption{Root fields of the forecastability contract, interface revision \texttt{1.0.0}. Types, units and structural constraints distinguish required members from optional metadata; nested records are listed separately.}\label{tab:contract-forecastability}\\
\toprule Field & JSON type & Unit & Root constraint & Required\\\midrule\endfirsthead
\toprule Field & JSON type & Unit & Root constraint & Required\\\midrule\endhead
\texttt{schema} & fixed value & metadata & const="chronosco\allowbreak{}pe.forecastabili\allowbreak{}ty/v1" & yes\\
\texttt{complete} & fixed value & metadata & const=true & yes\\
\texttt{protocol} & object & metadata &  & yes\\
\texttt{packages} & object & metadata &  & yes\\
\texttt{code\_\allowbreak{}sha256} & object & metadata &  & yes\\
\texttt{source\_\allowbreak{}sha256} & string & metadata & pattern="\textasciicircum{}[a-f0-9]\{64\}\$" & yes\\
\texttt{source} & object & metadata &  & yes\\
\texttt{source\_\allowbreak{}n} & fixed value & series & const=4227 & yes\\
\texttt{included\_\allowbreak{}n} & integer & counts & minimum=3000 & yes\\
\texttt{excluded} & array & metadata &  & yes\\
\texttt{model\_\allowbreak{}names} & array & metadata &  & yes\\
\texttt{summary} & object & mixed; see text &  & yes\\
\texttt{series} & array & mixed; see text & minItems=3000 & yes\\
\texttt{audit} & object & metadata &  & yes\\
\bottomrule
\end{longtable}
\endgroup


\subsection{Combination index: inventory, exclusions and compact summaries}
\label{sec:contract-combination-index}
The combination-index boundary is the compact navigation and summary layer for three cohorts: weekly, monthly and selected cases. Table~\ref{tab:contract-combination-index} describes interface revision \texttt{1.0.0}. Its completeness flag, experiment identity, protocol, implementation digests, package versions, source records and independent audit prevent a partial export from appearing to be the complete comparison. The case array has the expected 1,802 inventory entries, including identities excluded from scored inference. This differs from the 1,716 admitted external series and from the number of outer origins or prediction values. Browser forecast precision is fixed at eight decimal places, while the compressed numerical audit retains the underlying floating-point paths. The selected-case source binding prevents development examples from being joined to an unrelated reconstruction. Summaries average origins within series and then series within each external collection. They therefore require their stated denominators and cannot be reconstructed by giving every origin equal weight across an entire mixed collection. The browser first reads this index and fetches individual case documents on demand. A missing study is reported as unavailable; another experiment is not substituted under the same cohort labels or method names.
\begingroup\small
\begin{longtable}{@{}p{.22\textwidth}p{.14\textwidth}p{.14\textwidth}p{.29\textwidth}p{.09\textwidth}@{}}
\caption{Root fields of the combination-index contract, catalog version 1.0.0. Types, units and structural constraints distinguish required members from optional metadata; nested records follow the cited schema.}\label{tab:contract-combination-index}\\
\toprule Field & JSON type & Unit & Root constraint & Required\\\midrule\endfirsthead
\toprule Field & JSON type & Unit & Root constraint & Required\\\midrule\endhead
\texttt{schema} & fixed value & metadata & const="chronosco\allowbreak{}pe.combination-i\allowbreak{}ndex/v1" & yes\\
\texttt{complete} & fixed value & metadata & const=true & yes\\
\texttt{identity} & string & metadata & pattern="\textasciicircum{}[a-f0-9]\{64\}\$" & yes\\
\texttt{protocol} & object & metadata & Nested content follows the normative schema & yes\\
\texttt{code\_\allowbreak{}sha256} & object & metadata & Nested content follows the normative schema & yes\\
\texttt{export\_\allowbreak{}code\_\allowbreak{}sha256} & string & metadata & pattern="\textasciicircum{}[a-f0-9]\{64\}\$" & yes\\
\texttt{packages} & object & metadata & Nested content follows the normative schema & yes\\
\texttt{selected\_\allowbreak{}case\_\allowbreak{}source\_\allowbreak{}sha256} & string & metadata & pattern="\textasciicircum{}[a-f0-9]\{64\}\$" & yes\\
\texttt{cohorts} & array & metadata & minItems=3; maxItems=3 & yes\\
\texttt{cases} & array & metadata & minItems=1802; maxItems=1802 & yes\\
\texttt{sources} & object & metadata & Nested content follows the normative schema & yes\\
\texttt{browser\_\allowbreak{}forecast\_\allowbreak{}decimals} & fixed value & counts & const=8 & yes\\
\texttt{audit} & object & metadata & Nested content follows the normative schema & yes\\
\bottomrule
\end{longtable}
\endgroup


\subsection{Combination case: a source-bound origin and its allocation}
\label{sec:contract-combination-case}
The per-case combination boundary carries a source identity, dataset label, sequence length and digest, origin records and any explicit exclusion or representative relationship. Table~\ref{tab:contract-combination-case} describes interface revision \texttt{1.0.0}. An admitted origin retains its supplied prefix start and length, exact prefix digest, positive historical scale, reserved targets and method records. Each forecast row contains three quantile columns, while weights, selected configuration and validation metadata expose how a combination was obtained. Forecasts, targets and scale have source units; normalized scores and allocation weights are dimensionless. Origin indices and memory positions count observations. The loader restricts case paths to the declared cohort families and verifies the returned dataset and identifier. The walkthrough additionally requires the complete little-endian float64 prefix digest, horizon, seasonal lag and target agreement before showing weights beside an active comparison. Length and subject-name agreement alone are insufficient because two different histories can have those properties in common. An excluded case can legitimately have no scored origins; its absence from a leaderboard must remain traceable to its exclusion record. A missing requested document produces an explicit loading failure rather than a generated replacement sequence or weights borrowed from another case.
\begingroup\small
\begin{longtable}{@{}p{.22\textwidth}p{.14\textwidth}p{.14\textwidth}p{.29\textwidth}p{.09\textwidth}@{}}
\caption{Root fields of the combination-case contract, catalog version 1.0.0. Types, units and structural constraints distinguish required members from optional metadata; nested records follow the cited schema.}\label{tab:contract-combination-case}\\
\toprule Field & JSON type & Unit & Root constraint & Required\\\midrule\endfirsthead
\toprule Field & JSON type & Unit & Root constraint & Required\\\midrule\endhead
\texttt{schema} & fixed value & metadata & const="chronosco\allowbreak{}pe.combination-c\allowbreak{}ase/v1" & yes\\
\texttt{id} & string & metadata & No additional root-field range & yes\\
\texttt{dataset} & fixed value & metadata & enum=["m4\_\allowbreak{}weekly","m3\_\allowbreak{}monthly","selected\_\allowbreak{}cases"] & yes\\
\texttt{source\_\allowbreak{}n} & integer & observations & minimum=0 & yes\\
\texttt{sequence\_\allowbreak{}sha256} & string & metadata & pattern="\textasciicircum{}[a-f0-9]\{64\}\$" & yes\\
\texttt{origins} & array & metadata & Nested content follows the normative schema & yes\\
\texttt{excluded} & string & metadata & No additional root-field range & no\\
\texttt{representative} & string & metadata & No additional root-field range & no\\
\texttt{seasonality} & integer & observations & minimum=0 & no\\
\texttt{horizon} & integer & observations & minimum=0 & no\\
\bottomrule
\end{longtable}
\endgroup


\subsection{Release readiness and the judged tree}
\label{sec:contract-readiness}
Table~\ref{tab:contract-readiness} describes the offline readiness record. The catalog digest binds the normalized UTF-8 catalog text and the product version binds the checked revision. Contract entries follow the catalog order. Each entry identifies the contract revision and its declared status. A missing list contains the required document patterns for which the judged tree contains no file. Availability records the presence of required files; schema and scientific checks independently establish whether their contents are valid. The local stage is checked with the provider from the product checkout. Missing files, stale catalog or product bindings, and changed frozen research inputs stop admission.

Each positive glob must match at least one file. Negative globs limit the document sets considered by the schema checker and add no file requirement. A positive glob that expects an entire cohort can still match a partial cohort, so the separate scientific validator checks the complete identities, bindings and counts. Readiness alone does not attest the number of forecasts or cases.
\begingroup\small
\begin{longtable}{@{}p{.22\textwidth}p{.14\textwidth}p{.14\textwidth}p{.29\textwidth}p{.09\textwidth}@{}}
\caption{Root fields of the readiness contract, interface revision \texttt{1.0.0}. Types, units and structural constraints distinguish required members from optional metadata; nested records are listed separately.}\label{tab:contract-readiness}\\
\toprule Field & JSON type & Unit & Root constraint & Required\\\midrule\endfirsthead
\toprule Field & JSON type & Unit & Root constraint & Required\\\midrule\endhead
\texttt{schema} & fixed value & metadata & const="chronosco\allowbreak{}pe.readiness/v1" & yes\\
\texttt{product\_\allowbreak{}version} & string & metadata & pattern="\textasciicircum{}[0-9]+\allowbreak{}\textbackslash{}\textbackslash{}.[0-9]+\textbackslash{}\textbackslash{}.[0-9\allowbreak{}]+\$" & yes\\
\texttt{catalog\_\allowbreak{}sha256} & string & metadata & pattern="\textasciicircum{}[0-9a-f]\{64\}\$" & yes\\
\texttt{contracts} & array & metadata & minItems=1 & yes\\
\bottomrule
\end{longtable}
\endgroup

